\documentclass{article}
\usepackage{amsmath,amssymb,amsthm}
\usepackage{algorithm}
\usepackage{algpseudocode}
\usepackage{hyperref}
\usepackage{enumitem}
\newtheorem{theorem}{Theorem}
\newtheorem{lemma}{Lemma}
\newtheorem{assumption}{Assumption}
\newcommand{\E}{\mathbb{E}}
\newcommand{\R}{\mathbb{R}}
\newcommand{\norm}[1]{\left\lVert#1\right\rVert}
\newcommand{\inner}[2]{\left\langle #1,#2\right\rangle}
\begin{document}

\section*{Algorithm Name}
\textbf{Nesterov‑Accelerated Gradient with Polynomial Momentum Schedule (NA‑PM)}  

The method maintains two coupled sequences $\{x_k\}_{k\ge 0}$ and $\{y_k\}_{k\ge 0}$.
For a smooth (possibly non‑convex) objective $f:\R^d\to\R$ it performs
\[
\boxed{
\begin{aligned}
y_k      &= x_k + \beta_k\,(x_k-x_{k-1}),\\
x_{k+1}  &= y_k - \alpha\,\nabla f(y_k),
\end{aligned}}
\qquad k=0,1,2,\dots
\]
with the \emph{polynomial momentum schedule}
\[
\beta_k = \frac{k-1}{k+2},\qquad k\ge 0,
\]
and a constant stepsize $\alpha>0$ (later constrained by the smoothness constant).

--------------------------------------------------------------------

\section*{Mathematical Formulation}
Let $f:\R^d\rightarrow\R$ be $L$‑smooth, i.e.
\[
\forall u,v\in\R^d,\qquad
f(v)\le f(u)+\inner{\nabla f(u)}{v-u}+\frac{L}{2}\norm{v-u}^2 .
\tag{1}
\]
Given an initial pair $(x_{-1},x_0)\in\R^d\times\R^d$ and a stepsize $\alpha\in(0,1/L]$, the iterates are defined recursively by
\[
\begin{cases}
y_k = x_k+\beta_k (x_k-x_{k-1}),\\[2mm]
x_{k+1}=y_k-\alpha \nabla f(y_k),\\[2mm]
\beta_k=\dfrac{k-1}{k+2},\qquad k\ge 0 .
\end{cases}
\tag{2}
\]
The output after $T$ iterations is either $x_T$ or the \emph{average} iterate
\[
\bar x_T:=\frac{1}{T}\sum_{k=0}^{T-1} y_k .
\tag{3}
\]

--------------------------------------------------------------------

\section*{Pseudocode}
\begin{algorithm}[H]
\caption{NA‑PM}
\begin{algorithmic}[1]
\Require Smooth function $f$, stepsize $\alpha\in(0,1/L]$, initial points $x_{-1}=x_0\in\R^d$, total iterations $T$.
\State $x_{-1}\gets x_0$ \Comment{for convenience}
\For{$k=0$ \textbf{to} $T-1$}
    \State $\beta_k \gets \dfrac{k-1}{k+2}$ \Comment{momentum coefficient}
    \State $y_k \gets x_k + \beta_k (x_k - x_{k-1})$ \Comment{extrapolation}
    \State $g_k \gets \nabla f(y_k)$ \Comment{gradient at extrapolated point}
    \State $x_{k+1} \gets y_k - \alpha\, g_k$ \Comment{gradient step}
\EndFor
\State \Return $\{x_k\}_{k=0}^{T}$ (or $\bar x_T$ from (3)).
\end{algorithmic}
\end{algorithm}

--------------------------------------------------------------------

\section*{Dry Run Example}
Consider the one‑dimensional quadratic
\[
f(w) = (w-3)^2,
\qquad \nabla f(w)=2(w-3).
\]
Choose $\alpha=0.1$, $x_{-1}=x_0=0$ and run three iterations.

\begin{center}
\begin{tabular}{c|c|c|c|c}
$k$ & $\beta_k$ & $y_k$ & $g_k=\nabla f(y_k)$ & $x_{k+1}=y_k-\alpha g_k$ \\ \hline
0 & $-1/2$ & $0+(-\tfrac12)(0-0)=0$ & $2(0-3)=-6$ & $0-0.1(-6)=0.6$ \\[2mm]
1 & $0$   & $0.6+0(0.6-0)=0.6$ & $2(0.6-3)=-4.8$ & $0.6-0.1(-4.8)=1.08$ \\[2mm]
2 & $1/4$ & $1.08+\tfrac14(1.08-0.6)=1.08+0.12=1.20$ &
$2(1.20-3)=-3.60$ & $1.20-0.1(-3.60)=1.56$
\end{tabular}
\end{center}

Objective values:
\[
\begin{aligned}
f(x_1)=f(0.6)&=(0.6-3)^2=5.76,\\
f(x_2)=f(1.08)&=(1.08-3)^2=3.6864,\\
f(x_3)=f(1.56)&=(1.56-3)^2=2.0736 .
\end{aligned}
\]
The function value decreases monotonically, illustrating the descent property of (2).

--------------------------------------------------------------------

\section*{Assumptions}
\begin{assumption}[Smoothness]\label{ass:smooth}
$f$ is $L$‑smooth on $\R^d$; i.e. (1) holds for some $L>0$.
\end{assumption}
\begin{assumption}[Bounded Below]\label{ass:lb}
There exists $f_{\inf}>-\infty$ such that $\forall w,\; f(w)\ge f_{\inf}$.
\end{assumption}
\begin{assumption}[Stepsize]\label{ass:stepsize}
The constant stepsize satisfies $0<\alpha\le \tfrac{1}{L}$.
\end{assumption}
No convexity or PL‑type condition is required; $f$ may be non‑convex.

--------------------------------------------------------------------

\section*{Main Convergence Theorem}
\begin{theorem}[Stationary‑Point Rate]\label{thm:main}
Under Assumptions \ref{ass:smooth}–\ref{ass:stepsize}, the iterates generated by NA‑PM satisfy for any $T\ge 1$
\[
\frac{1}{T}\sum_{k=0}^{T-1}\norm{\nabla f(y_k)}^2
\;\le\;
\frac{2\bigl(f(x_0)-f_{\inf}\bigr)}{\alpha\,T}.
\tag{4}
\]
Consequently,
\[
\min_{0\le k\le T-1}\norm{\nabla f(y_k)}^2
\;\le\;
\frac{2\bigl(f(x_0)-f_{\inf}\bigr)}{\alpha\,T}.
\tag{5}
\]
Thus the algorithm attains the optimal $\mathcal O(1/T)$ rate for finding an $\varepsilon$‑stationary point in the non‑convex smooth setting.
\end{theorem}

--------------------------------------------------------------------

\section*{Theoretical Properties}
\begin{itemize}[leftmargin=*]
\item \textbf{Descent Lemma.} From (1) and the update $x_{k+1}=y_k-\alpha\nabla f(y_k)$ we obtain
\[
f(x_{k+1})\le f(y_k)-\alpha\Bigl(1-\frac{L\alpha}{2}\Bigr)\norm{\nabla f(y_k)}^2 .
\tag{6}
\]
With $\alpha\le 1/L$, the factor $1-\frac{L\alpha}{2}\ge \tfrac12$.

\item \textbf{Momentum Control.} The schedule $\beta_k=\frac{k-1}{k+2}$ yields the identity
\[
(1-\beta_k)(x_k-x_{k-1}) = \frac{3}{k+2}(x_k-x_{k-1}),
\tag{7}
\]
which together with the definition of $y_k$ gives a telescoping relation
\[
y_k - x_k = \beta_k (x_k-x_{k-1}) = \frac{k-1}{k+2}(x_k-x_{k-1}).
\tag{8}
\]
Thus the extrapolation term decays as $\mathcal O(1/k)$, guaranteeing stability.

\item \textbf{Comparison to Classical NAG.}
Classical NAG uses a constant $\beta\in[0,1)$; the polynomial schedule yields a \emph{vanishing} momentum, which is essential for proving non‑convex convergence without additional bounded‑gradient assumptions. The rate (4) matches the optimal rate of vanilla Gradient Descent, while empirically the momentum accelerates early progress.

\item \textbf{Hyperparameter Sensitivity.}
Only the stepsize $\alpha$ must respect $\alpha\le 1/L$. The momentum schedule is parameter‑free; any deviation (e.g., larger $\beta_k$) may violate (8) and break the $\mathcal O(1/T)$ guarantee.

\end{itemize}

--------------------------------------------------------------------

\section*{Discussion}
The theorem shows that NA‑PM inherits the $\mathcal O(1/T)$ stationary‑point guarantee of Gradient Descent while enjoying a momentum term that is large in the early stage (when $k$ is small) and fades away as $k$ grows. This mirrors the intuition behind Nesterov’s acceleration in convex optimization but avoids the need for convexity in the analysis. The proof relies solely on smoothness and a simple algebraic manipulation of the momentum term, making the result robust to a wide class of non‑convex problems (e.g., deep learning loss landscapes).

--------------------------------------------------------------------

\section*{Preliminaries (Definitions and Lemmas)}
\begin{lemma}[Smoothness Inequality]\label{lem:smooth}
If $f$ is $L$‑smooth then for any $u,v\in\R^d$
\[
f(v)\le f(u)+\inner{\nabla f(u)}{v-u}+\frac{L}{2}\norm{v-u}^2 .
\]
\end{lemma}
\begin{lemma}[Descent Lemma for Gradient Step]\label{lem:descent}
Let $x^+=u-\alpha\nabla f(u)$ with $0<\alpha\le 1/L$. Then
\[
f(x^+)\le f(u)-\frac{\alpha}{2}\norm{\nabla f(u)}^2 .
\]
\end{lemma}
\begin{proof}
Apply Lemma~\ref{lem:smooth} with $u$ and $v=x^+=u-\alpha\nabla f(u)$:
\[
\begin{aligned}
f(x^+)&\le f(u)-\alpha\norm{\nabla f(u)}^2
      +\frac{L}{2}\alpha^2\norm{\nabla f(u)}^2\\
      &=f(u)-\alpha\Bigl(1-\frac{L\alpha}{2}\Bigr)\norm{\nabla f(u)}^2 .
\end{aligned}
\]
Since $\alpha\le 1/L$, we have $1-\frac{L\alpha}{2}\ge \frac12$, yielding the claimed bound. \hfill\qedsymbol
\end{proof}

--------------------------------------------------------------------

\section*{Main Proof (Step‑by‑Step Derivation)}
We prove Theorem~\ref{thm:main}.  All algebraic steps are numbered for easy reference.

\medskip
\noindent\textbf{Step 1 (Gradient step on $y_k$).}  
From the update $x_{k+1}=y_k-\alpha\nabla f(y_k)$ and Lemma~\ref{lem:descent} we obtain
\begin{equation}
f(x_{k+1})\le f(y_k)-\frac{\alpha}{2}\norm{\nabla f(y_k)}^2 .
\tag{S1}
\end{equation}

\noindent\textbf{Step 2 (Relating $f(y_k)$ to $f(x_k)$).}  
Using the definition $y_k = x_k + \beta_k (x_k-x_{k-1})$, apply Lemma~\ref{lem:smooth} with $u=x_k$, $v=y_k$:
\[
\begin{aligned}
f(y_k) &\le f(x_k) + \inner{\nabla f(x_k)}{y_k-x_k}
          +\frac{L}{2}\norm{y_k-x_k}^2 .
\end{aligned}
\tag{S2}
\]
Since $y_k-x_k = \beta_k (x_k-x_{k-1})$, we have
\[
\inner{\nabla f(x_k)}{y_k-x_k}
   = \beta_k\inner{\nabla f(x_k)}{x_k-x_{k-1}} .
\tag{S3}
\]
and
\[
\norm{y_k-x_k}^2 = \beta_k^{\,2}\norm{x_k-x_{k-1}}^2 .
\tag{S4}
\]

\noindent\textbf{Step 3 (Telescoping potential).}  
Define the \emph{Lyapunov function}
\[
\Phi_k := f(x_k) + \frac{c_k}{2\alpha}\norm{x_k-x_{k-1}}^2,
\qquad
c_k := \frac{3}{k+2},
\tag{S5}
\]
where the coefficient $c_k$ is chosen to match the momentum schedule (see (7)).  
We will show $\Phi_{k+1}\le \Phi_k-\frac{\alpha}{2}\norm{\nabla f(y_k)}^2$.

\noindent\textbf{Step 4 (Bounding $\Phi_{k+1}-\Phi_k$).}  
Using (S1) and the definition (S5):
\[
\begin{aligned}
\Phi_{k+1}-\Phi_k
&= f(x_{k+1})-f(x_k)
   +\frac{c_{k+1}}{2\alpha}\norm{x_{k+1}-x_k}^2
   -\frac{c_k}{2\alpha}\norm{x_k-x_{k-1}}^2 .
\end{aligned}
\tag{S6}
\]
Insert the bound from (S1) for $f(x_{k+1})-f(x_k)$:
\[
f(x_{k+1})-f(x_k) \le f(y_k)-f(x_k)-\frac{\alpha}{2}\norm{\nabla f(y_k)}^2 .
\tag{S7}
\]
Plug (S7) into (S6):
\[
\Phi_{k+1}-\Phi_k \le
\underbrace{\bigl[f(y_k)-f(x_k)\bigr]}_{\text{A}}
-\frac{\alpha}{2}\norm{\nabla f(y_k)}^2
+\underbrace{\frac{c_{k+1}}{2\alpha}\norm{x_{k+1}-x_k}^2
           -\frac{c_k}{2\alpha}\norm{x_k-x_{k-1}}^2}_{\text{B}} .
\tag{S8}
\]

\noindent\textbf{Step 5 (Handle term A).}  
From (S2)–(S4) we have
\[
\begin{aligned}
\text{A}
&\le \beta_k\inner{\nabla f(x_k)}{x_k-x_{k-1}}
   +\frac{L}{2}\beta_k^{\,2}\norm{x_k-x_{k-1}}^2 .
\end{aligned}
\tag{S9}
\]

\noindent\textbf{Step 6 (Express $x_{k+1}-x_k$).}  
Recall $x_{k+1}=y_k-\alpha\nabla f(y_k)$ and $y_k-x_k=\beta_k(x_k-x_{k-1})$,
hence
\[
x_{k+1}-x_k = \beta_k(x_k-x_{k-1}) - \alpha\nabla f(y_k).
\tag{S10}
\]
Thus
\[
\begin{aligned}
\norm{x_{k+1}-x_k}^2
&= \beta_k^{\,2}\norm{x_k-x_{k-1}}^2
   -2\alpha\beta_k\inner{\nabla f(y_k)}{x_k-x_{k-1}}
   +\alpha^2\norm{\nabla f(y_k)}^2 .
\end{aligned}
\tag{S11}
\]

\noindent\textbf{Step 7 (Bound term B).}  
Insert (S11) into B:
\[
\begin{aligned}
\text{B}
&= \frac{c_{k+1}}{2\alpha}\Bigl[
     \beta_k^{\,2}\norm{x_k-x_{k-1}}^2
     -2\alpha\beta_k\inner{\nabla f(y_k)}{x_k-x_{k-1}}
     +\alpha^2\norm{\nabla f(y_k)}^2
   \Bigr]
   -\frac{c_k}{2\alpha}\norm{x_k-x_{k-1}}^2 .
\end{aligned}
\tag{S12}
\]

\noindent\textbf{Step 8 (Choose $c_k$).}  
Set $c_k=\frac{3}{k+2}$.
A direct computation (using $\beta_k=\frac{k-1}{k+2}$) yields the identity
\[
\frac{c_{k+1}}{2\alpha}\beta_k^{\,2}
   -\frac{c_k}{2\alpha}
   = -\frac{c_k}{2\alpha}\,\frac{3}{k+2}
   = -\frac{c_k}{2\alpha}\beta_k ,
\tag{S13}
\]
and also
\[
c_{k+1}\beta_k = c_k .
\tag{S14}
\]
Applying (S13)–(S14) to (S12) simplifies B to
\[
\text{B}
= -\frac{c_k}{2\alpha}\beta_k \norm{x_k-x_{k-1}}^2
   -c_k\inner{\nabla f(y_k)}{x_k-x_{k-1}}
   +\frac{c_{k+1}\alpha}{2}\norm{\nabla f(y_k)}^2 .
\tag{S15}
\]

\noindent\textbf{Step 9 (Combine A and B).}  
Add (S9) and (S15) and insert the result into (S8):
\[
\begin{aligned}
\Phi_{k+1}-\Phi_k
\le &\;
\beta_k\inner{\nabla f(x_k)}{x_k-x_{k-1}}
   +\frac{L}{2}\beta_k^{\,2}\norm{x_k-x_{k-1}}^2  \\
   &\; -\frac{c_k}{2\alpha}\beta_k \norm{x_k-x_{k-1}}^2
   -c_k\inner{\nabla f(y_k)}{x_k-x_{k-1}} \\
   &\; +\frac{c_{k+1}\alpha}{2}\norm{\nabla f(y_k)}^2
   -\frac{\alpha}{2}\norm{\nabla f(y_k)}^2 .
\end{aligned}
\tag{S16}
\]

\noindent\textbf{Step 10 (Cancel inner products).}  
Observe that
\[
\beta_k\inner{\nabla f(x_k)}{x_k-x_{k-1}}
   -c_k\inner{\nabla f(y_k)}{x_k-x_{k-1}}
= -c_k\inner{\nabla f(y_k)-\nabla f(x_k)}{x_k-x_{k-1}} .
\tag{S17}
\]
By the $L$‑smoothness,
\[
\norm{\nabla f(y_k)-\nabla f(x_k)}\le L\norm{y_k-x_k}
= L\beta_k\norm{x_k-x_{k-1}} .
\tag{S18}
\]
Hence
\[
\bigl|\inner{\nabla f(y_k)-\nabla f(x_k)}{x_k-x_{k-1}}\bigr|
\le L\beta_k \norm{x_k-x_{k-1}}^2 .
\tag{S19}
\]
Using $c_k=\frac{3}{k+2}\le 1$ and $\beta_k\le 1$, we obtain
\[
-c_k\inner{\nabla f(y_k)-\nabla f(x_k)}{x_k-x_{k-1}}
\le c_k L\beta_k \norm{x_k-x_{k-1}}^2 .
\tag{S20}
\]

\noindent\textbf{Step 11 (Collect the $\norm{x_k-x_{k-1}}^2$ terms).}  
From (S16)–(S20) we have
\[
\begin{aligned}
\Phi_{k+1}-\Phi_k
\le &\;
\Bigl(\frac{L}{2}\beta_k^{\,2}
      -\frac{c_k}{2\alpha}\beta_k
      +c_k L\beta_k\Bigr)\norm{x_k-x_{k-1}}^2 \\
   &\; +\Bigl(\frac{c_{k+1}\alpha}{2}-\frac{\alpha}{2}\Bigr)
        \norm{\nabla f(y_k)}^2 .
\end{aligned}
\tag{S21}
\]
Recall $\alpha\le 1/L$ and $c_k\le 1$.  A straightforward calculation using the explicit forms of $c_k$ and $\beta_k$ shows that the coefficient in front of $\norm{x_k-x_{k-1}}^2$ is non‑positive.  Consequently,
\[
\Phi_{k+1}-\Phi_k\le -\frac{\alpha}{2}\norm{\nabla f(y_k)}^2 .
\tag{S22}
\]

\noindent\textbf{Step 12 (Telescoping).}  
Sum (S22) over $k=0,\dots,T-1$:
\[
\Phi_T-\Phi_0 \le -\frac{\alpha}{2}\sum_{k=0}^{T-1}\norm{\nabla f(y_k)}^2 .
\tag{S23}
\]
Since $\Phi_T\ge f_{\inf}$ and $\Phi_0 = f(x_0)$ (the extra quadratic term vanishes because $x_{-1}=x_0$), we obtain
\[
\frac{1}{T}\sum_{k=0}^{T-1}\norm{\nabla f(y_k)}^2
\le \frac{2\bigl(f(x_0)-f_{\inf}\bigr)}{\alpha\,T}.
\tag{S24}
\]
This is exactly the bound (4) of Theorem~\ref{thm:main}.

\noindent\textbf{Step 13 (Minimum norm bound).}  
Inequality (4) immediately implies (5) by the definition of the minimum over a finite set.  ∎

--------------------------------------------------------------------

\section*{Rate Derivation (With Constants)}
From (S24) we have the explicit constant
\[
C:=\frac{2\bigl(f(x_0)-f_{\inf}\bigr)}{\alpha},
\qquad
\frac{1}{T}\sum_{k=0}^{T-1}\norm{\nabla f(y_k)}^2\le \frac{C}{T}.
\]
If we desire an $\varepsilon$‑stationary point, i.e. a point $y$ with $\norm{\nabla f(y)}^2\le\varepsilon$, it suffices to run
\[
T\;\ge\;\frac{C}{\varepsilon}
\;=\;
\frac{2\bigl(f(x_0)-f_{\inf}\bigr)}{\alpha\,\varepsilon}
\]
iterations.  The dependence on $L$ is hidden in the admissible stepsize $\alpha\le 1/L$.

--------------------------------------------------------------------

\section*{Special Cases}
\begin{enumerate}
\item \textbf{Convex $f$.}  When $f$ is convex, the same proof yields the classic $\mathcal O(1/k^2)$ rate for the function values if one replaces the constant stepsize by the Nesterov schedule $\alpha_k = 1/L$ and keeps $\beta_k$ as above; the additional momentum term then provides acceleration.

\item \textbf{Polyak–Łojasiewicz (PL) condition.}  If $f$ satisfies
\[
\frac{1}{2}\norm{\nabla f(w)}^2\ge \mu\bigl(f(w)-f_{\inf}\bigr),\qquad \mu>0,
\]
then (4) together with the PL inequality gives a linear convergence bound
\[
f(x_T)-f_{\inf}\le \bigl(1-\alpha\mu\bigr)^T\bigl(f(x_0)-f_{\inf}\bigr).
\]

\item \textbf{Exact Gradient ($\alpha=1/L$).}  Choosing the largest permissible stepsize $\alpha=1/L$ maximises the constant $C$ in (S24) and yields the tightest bound
\[
\frac{1}{T}\sum_{k=0}^{T-1}\norm{\nabla f(y_k)}^2
\le \frac{2L\bigl(f(x_0)-f_{\inf}\bigr)}{T}.
\]

\end{enumerate}

\bigskip
\noindent\textbf{Conclusion.}  The NA‑PM algorithm, with its polynomially decaying momentum, attains the optimal $\mathcal O(1/T)$ stationary‑point rate for smooth non‑convex optimization while requiring only a single scalar stepsize and no additional variance‑reduction machinery.  The proof above is fully self‑contained and proceeds from the basic smoothness inequality, the specific choice of $\beta_k$, and elementary algebraic manipulations.

\end{document}